% ============================================================
%  Table 3 — Component Ablation (SVD-LLMv2 Table-4 style)
%  LLaMA-7B, WikiText-2, 20%, 40%, and 60% compression ratios.
%  FR-SVD (F): Fisher-guided CEALC only
%  FR-SVD (R): Ordinary activation-aware CEALC with residual compensation
%  FR-SVD    : Both components
%  Required preamble packages:
%    \usepackage{booktabs}
%    \usepackage[table]{xcolor}
%    \usepackage{graphicx}   % for \resizebox
% ============================================================

\providecolor{mygreen}{rgb}{0.00,0.55,0.00}

\begin{table}[t]
\centering
\footnotesize
\renewcommand{\arraystretch}{1.15}
\setlength{\tabcolsep}{3pt}

\caption{Perplexity ($\downarrow$) of LLaMA-7B compressed by SAES-SVD and
FR-SVD with individual\,/\,both components under 20\%, 40\%, and 60\% compression ratios on
WikiText-2. FR-SVD (F) uses Fisher-guided CEALC truncation only; FR-SVD (R)
uses ordinary activation-aware CEALC followed by eigenspace residual
compensation; FR-SVD combines Fisher-guided CEALC with compensation.
For the complete FR-SVD model, relative perplexity reduction against
SAES-SVD is marked in \textcolor{mygreen}{green}.}
\label{tab:ablation-components}

\resizebox{\linewidth}{!}{%
\begin{tabular}{c|c|c|c|c}
\midrule
Ratio & SAES-SVD & FR-SVD (F) & FR-SVD (R) & FR-SVD \\ \midrule
0.2 & 7.17
 & 7.08
 & 7.12
 & \textbf{7.02}\,(\textcolor{mygreen}{$\downarrow$2.1\%}) \\
0.4 & 10.42
 & 10.02
 & 10.27
 & \textbf{9.81}\,(\textcolor{mygreen}{$\downarrow$5.9\%}) \\
0.6 & 22.01
 & 20.42
 & 21.13
 & \textbf{19.36}\,(\textcolor{mygreen}{$\downarrow$12.0\%}) \\ \midrule
\end{tabular}%
}
\end{table}
